\documentclass[UTF8,a4paper,11pt,fontset=fandol]{ctexart}
\usepackage[margin=22mm,headheight=16pt]{geometry}
\usepackage{amsmath,amssymb,newtxtext,newtxmath,booktabs,longtable,array,fancyhdr,xcolor,hyperref}
\hypersetup{unicode=true,colorlinks=true,urlcolor=blue!45!black,linkcolor=black,pdftitle={液态水散射本科生研究：全期规划与执行手册},pdfauthor={}}
\urlstyle{same}
\newcolumntype{P}[1]{>{\raggedright\arraybackslash}p{#1}}
\setlength{\parindent}{2em}\setlength{\parskip}{0.35em}\linespread{1.22}
\setlength{\emergencystretch}{3em}\widowpenalty=10000\clubpenalty=10000
\ctexset{section={format=\Large\bfseries,beforeskip=0.8ex,afterskip=0.8ex},subsection={format=\large\bfseries,beforeskip=1ex,afterskip=0.5ex}}
\pagestyle{fancy}\fancyhf{}\fancyhead[L]{\small 本科生液态水散射研究}\fancyhead[R]{\small 2026.10—2027.03}\fancyfoot[C]{\thepage}
\begin{document}
\begin{center}{\LARGE\bfseries 液态水散射本科生研究}\par\medskip{\Large 全期规划与执行手册}\par\medskip 2026年10月6日—2027年3月31日\end{center}
\section*{1　项目定位与启动任务书}
\subsection*{项目任务书}
暂定题目：液态水布里渊散射的谱线模拟、测量验证与误差分析。

\begin{longtable}{P{2.5cm}P{12.1cm}}
\toprule
\textbf{项目项} & \textbf{约定}\\\midrule\endhead
期限 & 2026年10月6日至2027年3月31日；正式结题日待导师确认。\\[0.35em]
团队与时间 & 2–4人，每人每周4–6小时；考试周和寒假减量。\\[0.35em]
当前条件 & 器材接近零；有导师，具体实验支持待落实；10万元为预算上限，非已批准额度。\\[0.35em]
主目标 & 完成水的布里渊频移及温度依赖验证，能解释标定、重复性和主要误差。\\[0.35em]
进阶与边界 & 线宽仅在响应条件具备后推进；瑞利组分用于理论和背景分析；热瑞利宽度、体积黏度与新机制非必达。\\[0.35em]
成功标准 & 文献→模型→标定或参考数据→验证→误差解释→报告；不以发表论文或发现新现象为前提。\\[0.35em]
当前状态 & 执行资料已准备；成员分工、经费、预约、采购、模拟运行和实验采集均待实际执行。\\[0.35em]
\bottomrule\end{longtable}
\subsection*{10月启动时需导师明确}
实测是否必需；借测能否用于结题；正式截止日期；实际资金与到账；可指导实验的人员；可借设备及允许导出的数据。未确认项记录在11号表，不替导师预填同意。

\subsection*{本手册如何使用}
先完成01号任务追踪表的T01–T04，再由导师确认技术路线。05号表只是采集计划；06–08号表是空白记录结构，只有实际取得数据后填写。阅读卡列出阅读目标，不代表精读已完成。

版本：v1.0，2026年10月6日。项目负责人：待分配；导师确认日期：待填写。已有两篇综述作为背景资料，本文作为独立执行手册，不覆盖原综述。

\clearpage
\section*{2　全期路线与转向条件}
\begin{longtable}{P{2.6cm}P{5.7cm}P{5.8cm}}
\toprule
\textbf{时间段} & \textbf{主要工作} & \textbf{阶段证据}\\\midrule\endhead
10/06–10/18 & 明确边界、分工、六篇精读、设备核查 & 任务书、设备表、阅读卡；要求仍未明确须登记\\[0.35em]
10/19–10/31 & 整理模型、参数与单位；核查借测及采购 & G1：导师确认路线和时间；不凭单件价格推断可行\\[0.35em]
11/01–11/15 & 统一数据处理；数值、噪声与失配验证 & G2：模型验证记录；无设备日期即启用备用渠道\\[0.35em]
11/16–12/15 & 指导操作、参考背景、室温水谱与排错 & G3：稳定测量可行性；不能只凭一次出现峰\\[0.35em]
12/16–01/10 & 锁定条件，三批次温度验证和误差分析 & 45条计划样品谱；参考背景另计；形成主图\\[0.35em]
01/11–02/28 & 离线整理、敏感性分析和报告初稿 & 可重跑数据包、初稿、按结论排序的补测清单\\[0.35em]
03/01–03/15 & 补参考与重复，必要时小型扩展 & G4：3月15日冻结数据和结论，不换主路线\\[0.35em]
03/16–03/31 & 交叉重跑、报告、PPT、模拟答辩 & 五项交付物；按正式截止日整体提前\\[0.35em]
\bottomrule\end{longtable}
\subsection*{四个决策节点}
10月31日：选择路线。11月15日：确认设备日期，否则并行寻找备用数据。12月15日：仍不能稳定取得水谱，停止扩大自建范围，由导师决定集中借测或调整任务。3月15日：冻结数据，后续以整理、复核和答辩为主。

若结题强制自测，仿真不能自行替代。设备条件不足时，应由导师协调借测或正式调整要求。任何路线变更填写09号会议表及11号变更表，保留决定与证据。

\clearpage
\section*{3　首两周行动与团队协作}
\subsection*{先把第一阶段做实}
\begin{longtable}{P{2.3cm}P{2.1cm}P{9.7cm}}
\toprule
\textbf{期限} & \textbf{角色} & \textbf{行动与完成证据}\\\midrule\endhead
10月8日前 & 全组 & 落实角色和固定讨论时段；每人写可投入时间与已有技能。\\[0.35em]
10月12日前 & B牵头 & 与导师讨论六项未决事项；记录确认结果，尚未回复的仍记待确认。\\[0.35em]
10月12日前 & A/C & 共同读两篇综述的核心公式、仪器响应和标定章节，各列三条不理解的问题。\\[0.35em]
10月18日前 & A牵头 & 完成六张阅读卡；每卡至少一个关键页码/图号、一项可用于本项目的结论、一项适用边界。\\[0.35em]
10月18日前 & B牵头 & 按功能核实设备和借测条件，区分能用、需预约、需采购和未知。\\[0.35em]
10月18日前 & C牵头 & 建立原始数据、处理结果、图表与文档目录；试用CSV模板，不填造实验数据。\\[0.35em]
\bottomrule\end{longtable}
\subsection*{分工与工作量}
A：理论与文献；B：实验与条件；C：数据与成果；D：交叉复核（可选）。两人团队分为A+C、B；三人团队按A/B/C分工并交叉复核；四人团队增加D。真实姓名在启动讨论时填写。

每周每人约2小时学习或理论、2小时实际任务、0.5–1小时记录与复核；实验周可集中安排，不能把名义投入时间等同于设备可用时间。每周30分钟小组讨论，每两周一次导师讨论。

\subsection*{每次汇报四句话}
本周期完成了什么；证据在哪个文件；目前具体卡在哪里；下一步由谁在何时完成。以结果图、原始记录或决策文件为证据，不只汇报阅读时长。

\subsection*{两人至四人的复核要求}
即使只有两人，采集或分析也要由另一人核查单位、样品编号和结论。最后的重跑由未主导处理流程的成员完成；导师指导与成员复核分别记录。

\clearpage
\section*{4　模型验证的最低完成标准}
\subsection*{模型的输入与输出}
输入包括真空波长、水折射率、声速、内部散射角、样品状态、谱线宽度与面积，以及仪器响应、采样和噪声。每项数值标为实测、厂家规格或假设，并写来源；不把临时示例当作未来器材参数。

\begin{equation*}
q=\frac{4\pi n_w}{\lambda_0}\sin\frac{\theta_w}{2},\qquad
\nu_B=\frac{c_s q}{2\pi}.
\end{equation*}
水中散射角决定样品波矢；FP腔内角和腔介质决定分析器响应。两者分开。输出至少包括合成数据、拟合参数、偏差、重复离散与失败原因；成像FP还应保留空间坐标到频率的映射。

\begin{longtable}{P{2.0cm}P{7.3cm}P{4.8cm}}
\toprule
\textbf{检查} & \textbf{具体执行} & \textbf{保留证据}\\\midrule\endhead
公式与无噪声 & 独立重算频移；检查Hz/角频率及FWHM/HWHM；恢复已知合成参数 & 输入真值、恢复值、数值偏差\\[0.35em]
收敛 & 同一物理问题采用原步长、半步长、四分之一步长 & 参数/图像差异；预先锁定数值误差预算\\[0.35em]
噪声 & 至少3个水平，每个50次独立实现；同一真值，记录噪声定义和种子 & 150次原始结果、均值偏差、标准差、失败率\\[0.35em]
失配 & 分别改变生成数据的响应宽度、背景、几何；拟合仍用原模型 & 误差随失配改变的图，不能只看残差\\[0.35em]
外部依据 & 用独立公式、文献数据或自测结果补充模型自洽检查 & 区分理论验证、数据复现与实验验证\\[0.35em]
\bottomrule\end{longtable}
\subsection*{实施前锁定，运行后报告}
11月5日前在10号表对应配置中确定噪声类型与三档强度、种子规则、参数约束、拟合失败判据和数值预算。50次指各水平独立实现，不能把同一条数据重复拟合50次冒充独立噪声检验。失败样本计入分母并保留日志。

本手册已提供验证计划表，未声称模型程序或150次运行已完成。具体脚本与实际装置适配按11月阶段执行。

\clearpage
\section*{5　实验主线与采集操作框架}
\subsection*{采集设计：三个跨日批次，共45条样品谱}
同一来源纯水；每批进行一次独立准备或重装并记录，安排在不同日期。每批三个温度点，每点连续五条谱。参考、暗背景、空池等另计，不能占用45条样品谱编号。

\begin{longtable}{P{1.3cm}P{5.0cm}P{2.0cm}P{5.3cm}}
\toprule
\textbf{批次} & \textbf{建议温度顺序} & \textbf{样品谱数} & \textbf{目的}\\\midrule\endhead
D1 & 20 → 25 → 30 ℃ & 15 & 建立首批结果\\[0.35em]
D2 & 25 → 30 → 20 ℃ & 15 & 减少所有批次都单向升温的顺序混杂\\[0.35em]
D3 & 30 → 20 → 25 ℃ & 15 & 进一步比较跨日重复与顺序影响\\[0.35em]
\bottomrule\end{longtable}
温度顺序是排程建议，不代表已经控制热历史；每次转温都应等待预先规定的稳定条件，并记录等待时间和实际温度。设备不适合上述顺序或温区时，采集前统一更改计划并说明原因。

\subsection*{一次批次的顺序}
① 记录样品与配置；② 在指导下检查功率、饱和及原始数据导出；③ 采集频率参考、暗背景和必要空池；④ 达到稳定条件后依计划采样，每个温度记录连续5条；⑤ 批次结束复测参考；⑥ 当天核查文件是否可读、ID是否匹配，再备份原始数据。

\subsection*{先标定，再划分验证数据}
使用独立频率参考时记录参考值及不确定度。若以水的已知频移标定，必须预先区分标定状态与留作检验的状态，并在报告中说明共享的物性输入；不能把同一状态同时称为标定依据和独立验证。

\subsection*{统计与质量判据}
单日同温度5条用于短期重复性；同温度跨3日用于批次差异。样本量较小，结论限于这三批测量，不能宣称长期稳定性。温度稳定阈值、观察窗口、漂移处理和精度目标在预实验后锁定，不预先承诺任意百分比准确度。

饱和、断采、温度未稳或参考丢失等异常必须保留原文件并标记；补测使用新ID。05号计划表与06号实际记录分别填写，计划完成数和有效测量数分别统计。

\clearpage
\section*{6　设备路线、经费与备用方案}
\subsection*{先获取完整测量能力，再决定购买什么}
优先借用或预约成熟系统；其次在成熟光路上补模块；最后才考虑从零搭建。当前器材接近零，不以10万元上限推定一定能买齐。必须有可指导的人、能导出的数据、可执行的参考方案以及明确的设备可用日期。

\begin{longtable}{P{2.4cm}P{8.0cm}P{3.7cm}}
\toprule
\textbf{候选路线} & \textbf{10月底前必须得到的证据} & \textbf{不满足时}\\\midrule\endhead
借测/共享 & 可测水样、预约日期、参与方式、数据使用与导出范围、完整费用 & 继续核实其他渠道，同时推进模型\\[0.35em]
局部补充 & 已有主光路可工作；缺项及接口明确；报价、到货和调试人确定 & 不把缺项称为完整系统\\[0.35em]
从零搭建 & 全链器件与兼容、总价、交期、标定方法、指导与排错时间 & 不进入主进度，不进行零散大件采购\\[0.35em]
\bottomrule\end{longtable}
\subsection*{经费控制}
以实际批准总额B为准，规划约0.2B预留、约0.8B分阶段安排。若B最终为10万元，则预留约2万元；这只是预算规则，不是厂家报价。设备款、税费运输、必要附件、控制软件、测试与维修应按实际交付范围核实。

\subsection*{备用路线落实}
11月15日无可靠设备日期：继续模型验证，同时寻找合法可用的原始水谱、补充数据或作者数据。记录许可、原始链接、获取日期和使用条件。未取得原始数据时，读图只作趋势与量级比较，保存数字化误差与数据类别。

12月15日仍无稳定水谱：停止扩大自建范围，向导师汇报已完成工作、阻碍和替代方案。若要求自测，需协调集中借测或正式修改任务书；不得未经确认降为纯仿真结题。

\subsection*{可选扩展的进入条件}
只有主线数据、参考、误差分析和报告初稿已经完整，才开展积分时间对照。记录总时间和漂移，不只比较单次拟合误差。若占用补齐主证据的时间，直接取消并登记不适用。

\clearpage
\section*{7　数据管理与证据复核}
\subsection*{文件组织与数据边界}
执行资料包提供数据工作区：00\_配置、01\_原始数据、02\_参考背景、03\_处理结果、04\_图表、05\_报告、06\_会议。原始数据只追加，不原位修改；每个处理版本另存并记录代码和配置版本。

\begin{longtable}{P{1.7cm}P{5.1cm}P{7.3cm}}
\toprule
\textbf{表号} & \textbf{用途} & \textbf{填写原则}\\\midrule\endhead
01–04 & 任务、设备、预算和参数 & 状态以证据为准；空报价不填0元，未知参数不冒充实测\\[0.35em]
05 & 45条计划样品谱 & 用于安排采集，实际日期和谱ID采集后填写\\[0.35em]
06–07 & 样品与参考背景原始记录 & 通过唯一ID、时间和配置关联，不用文件名猜测条件\\[0.35em]
08 & 拟合与验证结果 & 模拟、文献、实测分开；失效和剔除理由保留\\[0.35em]
09–11 & 会议、模型检查和变更 & 决定有负责人和日期；变更保留前提与原因\\[0.35em]
\bottomrule\end{longtable}
\subsection*{不确定度与比较}
短期重复性、跨日变化、频率标定、温度、几何和响应均需评估。统计标准差不等于所有误差；未定量的来源单独说明，不能默认为零。与参考关系比较时，共享温度、折射率或标定数据会带来相关性。

\begin{equation*}
\Delta=\widehat{\nu}_B-\nu_{B,\mathrm{ref}},\qquad
u^2(\Delta)=u^2(\widehat{\nu}_B)+u^2(\nu_{B,\mathrm{ref}})
-2\operatorname{cov}(\widehat{\nu}_B,\nu_{B,\mathrm{ref}}).
\end{equation*}
这里u表示标准不确定度。仅在有依据认为独立时，才省略协方差。误差条和区间写明定义；“在区间内相符”说明相容程度，不代表已证明模型唯一正确。

\subsection*{三月交叉复核}
3月15日列出冻结的数据、代码、配置和图表版本。非主分析成员依据说明，从原始输入生成核心图表，核对数字与报告一致。任何补丁保留修改记录，不能只手工改报告中的数值。

\clearpage
\section*{8　结题成果与验收口径}
\subsection*{五项最终交付物}
\begin{longtable}{P{2.1cm}P{7.0cm}P{5.0cm}}
\toprule
\textbf{交付物} & \textbf{建议内容} & \textbf{验收方式}\\\midrule\endhead
研究报告 & 20–30页，项目要求优先；背景引用已有两篇综述，突出本组工作 & 每个主要结论有数据或来源；局限明确\\[0.35em]
分析包 & 原始输入、配置、程序、环境与运行说明 & 另一成员能够重跑核心结果\\[0.35em]
核心图表 & 预测、代表性谱、温度验证、重复性、误差或敏感性 & 图与原始记录对应；标明数据类别\\[0.35em]
流程说明 & 采集/验证步骤、状态条件、异常和失败处理 & 后续同学能够理解并接续\\[0.35em]
答辩PPT & 12–15页，说明问题、方法、证据与边界 & 模拟答辩后保留问题和修订记录\\[0.35em]
\bottomrule\end{longtable}
\subsection*{报告提纲与写作安排}
1 研究问题与范围；2 文献依据和可比较条件；3 模型、装置或数据来源；4 标定与数据处理；5 验证结果；6 重复性、不确定度和偏差解释；7 结论、局限与后续建议。附录放参数表、详细流程和数据索引。没有完成的实验不写成结果章节。

\subsection*{结题自查}
每张图能追溯到输入文件；每个拟合参数标单位与宽度定义；原始数据与修改记录完整；文献图读值和原始数据不混称；标定和验证不循环；模拟运行不冒充实测；负结果和失败有记录；未定的误差来源没有被忽略。

\subsection*{什么算有效成果}
可靠复现是成果；可重复的偏差及其定量定位也是成果。没有出现预期差异时，应说明在当前噪声、标定与样本量下能排除什么，不能断言物理效应不存在。没有独立标定时，结论限定为依赖当前参考的验证。

\subsection*{时间不足时的删减顺序}
先取消积分时间扩展，再取消线宽进阶，保留主频移验证、参考记录、误差说明与报告。如果主实验条件仍不成立，执行已由导师确认的备用路线，完整呈现该路线实际完成的证据。

\clearpage
\section*{9　核心文献阅读卡 1–3}
\subsection*{R1 Kabakova等，2024}
Brillouin microscopy

\noindent\href{https://doi.org/10.1038/s43586-023-00286-z}{DOI: 10.1038/s43586-023-00286-z}\par\smallskip
获取状态：方法综述；已取得正式全文PDF。

阅读重点：先读Experimentation、Results、Reproducibility与Limitations；建立测量链概念。

应交出的阅读成果：画出光源、样品、分析器、探测和标定的功能图；区分频移、线宽和纵向模量。

解释边界：论文涉及多种显微与生物应用，不能据此认为本项目需要全部功能。

负责人／完成日：待填；关键页码或图号：待填；自己的推导或复现：待填；未理解问题与复核意见：待填。

\subsection*{R2 Mountain，1966}
Thermal relaxation and Brillouin scattering in liquids

\noindent\href{https://doi.org/10.6028/jres.070A.017}{DOI: 10.6028/jres.070A.017}\par\smallskip
获取状态：理论原文；NIST全文PDF。

阅读重点：结合物理综述3.1–3.4读模型假设、热弛豫与中央组分。

应交出的阅读成果：制作Hz/角频率、HWHM/FWHM和衰减率对照表，解释为何线宽可包含多种耗散。

解释边界：不要求第一轮推导所有公式；先保证本项目实际使用部分的定义一致。

负责人／完成日：待填；关键页码或图号：待填；自己的推导或复现：待填；未理解问题与复核意见：待填。

\subsection*{R3 Koestel与Walther，2024}
The Brillouin linewidth in water as a function of temperature and salinity: the missing empirical relationship

\noindent\href{https://doi.org/10.1007/s00340-024-08189-x}{DOI: 10.1007/s00340-024-08189-x}\par\smallskip
获取状态：直接水样研究；作者仓储全文PDF。

阅读重点：优先读样品制备、530 nm/179°几何、频率轴与响应处理。

应交出的阅读成果：摘录温度、盐组成、参考方式与有效精细度的区别；画一张“参考输入/独立输出”表。

解释边界：摘要1.8–35 ℃与方法最低271.4 K不一致；海盐混合物不是纯NaCl；不能反用水标定结果独立验证同一关系。

负责人／完成日：待填；关键页码或图号：待填；自己的推导或复现：待填；未理解问题与复核意见：待填。

\clearpage
\section*{10　核心文献阅读卡 4–6}
\subsection*{R4 Yan等，2020}
Evaluation of commercial virtually imaged phase array and Fabry-Pérot based Brillouin spectrometers for applications to biology

\noindent\href{https://doi.org/10.1364/BOE.401087}{DOI: 10.1364/BOE.401087}\par\smallskip
获取状态：仪器比较；作者仓储全文PDF。

阅读重点：重点读纯水/盐水比较、共同测量条件及响应修正。

应交出的阅读成果：整理同条件比较所需的波长、几何、积分、响应和误差字段，并与设备表对应。

解释边界：旧型号和论文配置的性能不等于当前商品保证；不要直接照搬速度排名。

负责人／完成日：待填；关键页码或图号：待填；自己的推导或复现：待填；未理解问题与复核意见：待填。

\subsection*{R5 Bouvet等，2025}
Consensus statement on Brillouin light scattering microscopy of biological materials

\noindent\href{https://doi.org/10.1038/s41566-025-01681-6}{DOI: 10.1038/s41566-025-01681-6}\par\smallskip
获取状态：测量共识；出版社全文网页及补充资料。

阅读重点：优先读数据报告、频率/线宽标定、参考样品和不确定度相关内容。

应交出的阅读成果：把本项目必须记录的信息压缩为一页清单，解释重复性与准确度的区别。

解释边界：共识主要面向生物材料；本项目选用通用测量要求，不照搬所有生物成像要求。

负责人／完成日：待填；关键页码或图号：待填；自己的推导或复现：待填；未理解问题与复核意见：待填。

\subsection*{R6 Testi等，2026}
Stabilized real-time Brillouin microscopy reveals fractal organization of protein condensates in living cells

\noindent\href{https://doi.org/10.1038/s41467-026-68984-2}{DOI: 10.1038/s41467-026-68984-2}\par\smallskip
获取状态：近期系统论文；出版社全文PDF。

阅读重点：只重点阅读水的温度验证、实测响应与稳定性分析部分。

应交出的阅读成果：提取一组水验证图的输入、输出、参考与误差定义；说明自己能复现哪一级证据。

解释边界：生物凝聚体结果不属于本项目研究目标；取得作者数据前不把图像读值称为原始数据。

负责人／完成日：待填；关键页码或图号：待填；自己的推导或复现：待填；未理解问题与复核意见：待填。

\end{document}